\chapter{Criterios de revisión de conversaciones}
\label{app:instrumentos-revision}

Este apéndice desarrolla la escala del capítulo~\ref{chap:metodologia-evaluacion}.
Los dos jueces del jurado y la revisión humana aplican los mismos anclajes, que el
jurado recibe literalmente en sus instrucciones (versión \texttt{judge-v1}), y
ven la misma información del episodio (apartado~\ref{sec:metodo-mixta}).
Los ejemplos son ilustrativos, no anotaciones de episodios ejecutados.

\section{Dimensiones del asistente}

Se puntúa cada dimensión por separado y se indica la evidencia que justifica la
nota. El jurado da los números de los mensajes y copia un fragmento literal
breve de uno de ellos, que el programa comprueba. En la revisión humana se
anotan los mensajes que justifican la nota. La Tabla~\ref{tab:anclajes-asistente} concreta los
niveles. No se compensa un error central con respuestas correctas en otros
turnos.

\begingroup
\setlength{\tabcolsep}{4pt}
\small
\begin{longtable}{L{0.15\textwidth}L{0.18\textwidth}L{0.18\textwidth}L{0.18\textwidth}L{0.18\textwidth}}
\caption{Anclajes de la escala del asistente.}\label{tab:anclajes-asistente}\\
\toprule
\textbf{Dimensión} & \textbf{0} & \textbf{1} & \textbf{2} & \textbf{3} \\
\midrule
\endfirsthead
\toprule
\textbf{Dimensión} & \textbf{0} & \textbf{1} & \textbf{2} & \textbf{3} \\
\midrule
\endhead
Respuesta administrativa & Responde mal al objetivo central o anuncia un efecto falso. & Orientación insuficiente con un defecto importante para resolver la gestión. & Resuelve la necesidad con un defecto relevante, no decisivo. & Respuesta correcta y adecuada, incluido un rechazo justificado. \\
Continuidad & Pierde la tarea o ignora una corrección expresa. & Conserva fragmentos, pero obliga a reconstruir información importante. & Mantiene la tarea con alguna aclaración redundante. & Conserva información compatible y actualiza correcciones sin arrastre indebido. \\
Claridad & No se comprende qué significa la respuesta o qué hacer. & Ambigüedad importante en instrucciones o explicación. & Comprensible, pero requiere releer o aclarar un término. & Directa y comprensible, con la siguiente acción clara cuando corresponde. \\
Ausencia de repetición & Un bucle impide avanzar en la gestión. & Repite solicitudes o información de forma que dificulta avanzar. & Redundancia evitable, pero la gestión progresa. & No repite sin necesidad. Confirmaciones obligatorias no cuentan como redundancia. \\
Adecuación al perfil & Contradice una necesidad expresa o impone un estereotipo. & Omite una adaptación relevante y dificulta la interacción. & Adaptación parcial con un defecto no decisivo. & Usa necesidades y preferencias relevantes sin añadir información innecesaria. \\
Fidelidad material & Dato central falso o anuncio de una escritura inexistente. & Omisión o discrepancia importante respecto a los hechos disponibles. & Hechos centrales correctos con un defecto no decisivo de presentación. & Petición, datos comunicados y efectos observados concuerdan. \\
\bottomrule
\end{longtable}
\endgroup

Si Ana pide otra hora y el asistente sigue ofreciendo la que rechazó, existe
un problema de continuidad aunque el mensaje sea claro. Si el asistente anuncia
una reserva y la evidencia demuestra que no se creó, respuesta administrativa
y fidelidad material reciben 0. Claridad se valora según la comprensión del
texto, no se reduce automáticamente por ese error. Una negativa correcta por
falta de permiso puede recibir 3 en respuesta administrativa: no reservar no
equivale siempre a fallar.

\section{Dimensiones del usuario simulado}

Se juzgan sus mensajes frente al objetivo, la persona funcional y la
información pública permitida. Un error del asistente o de la agenda no se
atribuye al simulador. Se usa la misma graduación: 0 incumple la dimensión,
1 presenta un defecto importante, 2 es utilizable con un defecto no decisivo
y 3 conserva la conducta esperada. La Tabla~\ref{tab:anclajes-simulador}
especifica qué se observa en cada caso.

\begin{longtable}{L{0.23\textwidth}L{0.32\textwidth}L{0.33\textwidth}}
\caption{Conductas revisadas del usuario simulado.}\label{tab:anclajes-simulador}\\
\toprule
\textbf{Dimensión} & \textbf{Conducta adecuada} & \textbf{Ejemplo de incumplimiento} \\
\midrule
\endfirsthead
\toprule
\textbf{Dimensión} & \textbf{Conducta adecuada} & \textbf{Ejemplo de incumplimiento} \\
\midrule
\endhead
Fidelidad al objetivo & Mantiene la gestión o comunica un cambio permitido. & Acepta una cita incompatible con una condición que no ha retirado. \\
Coherencia & Sus mensajes y correcciones son compatibles. & Se contradice sobre una restricción sin reconocer el cambio. \\
Información permitida & Usa solo hechos que puede conocer y propuestas visibles. & Cita una operación o identificador oculto en el estado interno. \\
Ausencia de invención & No añade circunstancias que alteran el caso. & Inventa una cita previa o una autorización para continuar. \\
Confirmaciones correctas & Acepta, rechaza o aplaza según propuesta e intención expresas. & Confirma un recurso distinto del que la conversación permite seleccionar. \\
\bottomrule
\end{longtable}

Una corrección expresa no es incoherencia por sí sola. Tampoco se penaliza que
el usuario rehúse una operación permitida. En cada nota se explica si el defecto
impide cumplir la dimensión, la afecta de forma importante o deja una conducta
utilizable. El éxito del asistente no sustituye esa revisión.

\section{Aplicabilidad y evidencia}

\texttt{NA} significa que el caso no ofrece oportunidad para observar una
dimensión. Nunca se transforma en 0 ni en 3. Una sola respuesta sin intercambio
previo no permite evaluar continuidad. Una consulta de lectura seguida de un
agradecimiento no ofrece una confirmación que puntuar. En cambio, \enquote{no quiero
cancelar} permite revisar si el simulador rechaza correctamente la operación.

En cada pasada, cada episodio se revisa una sola vez, con la conversación y los
efectos observados en la agenda a la vista. Anunciar \enquote{he reservado} no demuestra la
escritura: el anuncio se contrasta con esos efectos y, si la reserva no
aparece, respuesta administrativa y fidelidad material reciben 0. Si un juez
no devuelve una respuesta válida, el episodio se registra como incidente del
jurado y no recibe ninguna nota.

El éxito material de la gestión no se decide con esta escala. Lo determina el
oráculo determinista a partir del estado final de la agenda
(capítulo~\ref{chap:metodologia-evaluacion}), y ninguna nota de 0 a 3 lo
modifica.
